% Generated from InterviewRecord_Group3.csv (one-off conversion; the CSV is final).
\begin{longtable}{@{}L{1.3cm} p{\dimexpr\linewidth-1.3cm-2\tabcolsep\relax}@{}}
\toprule
\textbf{ID} & \textbf{Notes} \\
\midrule
\endhead
\bottomrule
\endlastfoot
\multicolumn{2}{@{}p{\linewidth}@{}}{\textbf{G1.} What is your role in the club, and how many years have you been involved?} \\*
S3-P1 & Executive (Co-chair), Coach (Logistics), 3 years \\
S3-P2 & General member, participant, 5 years \\
S3-P3 & Executive (Finance Director), Team Captain, 6 years \\
\addlinespace
\multicolumn{2}{@{}p{\linewidth}@{}}{\textbf{G2.} Which tools do your teams use to design, analyze, and store their bridges?} \\*
S3-P1 & AutoCAD and Inventor for design, SAP model and hand calculations for analysis, all information stored on team members' personal laptops and MS Teams \\
S3-P2 & Autodesk Inventor to design, create engineering drawings, conduct FEA and validate rules compliance; Excel sometimes used to determine optimal designs \\
S3-P3 & Autodesk Inventor used to model the bridge in CAD over the last 6 years; analysis mostly hand calculations given the simplicity of the bridge; an Autodesk Inventor FEA model also used, but did not yield good results given the complexity of the material \\
\addlinespace
\multicolumn{2}{@{}p{\linewidth}@{}}{\textbf{G3.} What were the largest challenges you faced when you first joined the club?} \\*
S3-P1 & Figuring out where to start, no large archive/database for comparing past designs, strategies, the build/design process, and competition results \\
S3-P2 & Lack of collected information about past designs on joining in 2021; a few older members remembered how things had been done, but no documentation, photo archive, or technique guides; made it challenging to learn from others' mistakes and slowed progress as a team \\
S3-P3 & Figuring out what to do approaching the competition, a big disconnect between club leaders and the teams; second, space to build the bridge, with no place available at the time (team ended up using a basement); team did not really know what it was doing, a good starting point noted as very valuable \\
\addlinespace
\multicolumn{2}{@{}p{\linewidth}@{}}{\textbf{E1.} What information regarding the club administration is most important to consolidate?} \\*
S3-P1 & Past bridge data: overall design, calculations/analysis performed, competition results, successes/failures, etc. \\
S3-P2 & \textit{No answer recorded} \\
S3-P3 & Standardization of analysis and consolidation of all knowledge gained from building bridges, seen as leading to rapid improvement of the team over time \\
\addlinespace
\multicolumn{2}{@{}p{\linewidth}@{}}{\textbf{E2.} When you first became an executive or coach, what were your largest difficulties?} \\*
S3-P1 & Previously Finance executive, few difficulties there as the club's finances are straightforward; new co-chair position brought many difficulties: planning how to familiarize new members with how Troitsky works, identifying the most feasible bridge types (to keep students from being discouraged if their bridges do not perform well), getting information to members other than through in-person meetings \\
S3-P2 & \textit{No answer recorded} \\
S3-P3 & As co-president, the greatest difficulty was keeping all teams on time and on track; most teams did not know what they were doing or what they were getting into when they started \\
\addlinespace
\multicolumn{2}{@{}p{\linewidth}@{}}{\textbf{E3.} How would you like to see the software sub team integrated into the club administration?} \\*
S3-P1 & Dedicated team working closely with the co-chairs, executive team and coaches, similar to existing teams; designated lead who communicates regularly with the co-chairs, gives development updates, and works with other areas of the club to find improvements; long term, expand to support both technical and administrative sides of the competition; technical: interactive design tools where students input a bridge configuration, select design constraints, and get help with structural calculations, member forces and other engineering considerations (especially for students new to structural design); administrative: centralized team-management platform (milestone submission, deadline tracking, coach and executive feedback, progress monitoring) to streamline communication and reduce the number of separate platforms; since the executive team changes yearly, clear documentation, responsibilities and transition procedures needed so the Toolbox can keep being maintained; wants the software sub-team to be a permanent part of the club structure and the Toolbox a long-term resource \\
S3-P2 & \textit{No answer recorded} \\
S3-P3 & Consistent maintenance and improvement year after year; software sub-team to be a consistent part of the executive team \\
\addlinespace
\multicolumn{2}{@{}p{\linewidth}@{}}{\textbf{C1.} What tool have you found is the most accurate and useful in predicting your bridge's performance?} \\*
S3-P1 & Team's performance estimates not very accurate due to little material data for the composite material; SAP used to identify large weaknesses, tests run in SAP along with hand calculations to estimate the anticipated max load \\
S3-P2 & Identifies a weak point and calculates the expected stresses at that point using hand calculations and Excel; without material testing, knowledge of exact deformation under load and manufacturing quality, predicting performance accurately is very difficult \\
S3-P3 & First two years relied purely on intuition; last year, bridge designed from static calculations and formulas and optimized for weight using an Excel sheet, but quite inaccurate because the statics equations did not match the actual static system; no tool found extremely helpful; interest in implementing FEA in design choices \\
\addlinespace
\multicolumn{2}{@{}p{\linewidth}@{}}{\textbf{C2.} How close to a model estimate do you expect your bridge's performance to be, including uncertainties?} \\*
S3-P1 & Expects performance reasonably close to the model estimate, accounting for uncertainty in material properties, joint strength and construction quality; suggested the Toolbox give an estimated performance range rather than a single value, so teams can compare designs, identify weaknesses and understand potential failure points \\
S3-P2 & Aims to be within 20\%, though precise estimates have not been a focus of their teams in the past; better understanding and models noted as helpful; uncertainties play a very big role, especially in the material data used (mostly ultimate tensile/compressive strength) \\
S3-P3 & Depends on the model; 2-D truss models good for finding a theoretical maximum, given a good estimate of material failure; 3-D effects make building models tricky; uncertainties there need to be explored more, as a really good model is not possible without looking at them \\
\addlinespace
\multicolumn{2}{@{}p{\linewidth}@{}}{\textbf{C3.} What are the most common reasons for bridge failures in your experience?} \\*
S3-P1 & Most commonly torsion; first year, disqualified on reaching max deflection so the failure was never seen; other years, bridge loaded asymmetrically (placement correct, manufacturing issue) so only one truss was loaded, causing it to twist and break \\
S3-P2 & Misalignment of built parts (manufacturing quality) often causes bridges to twist or tip prematurely; weak joints between parts second most common \\
S3-P3 & Manufacturing defects most common; a poorly made bridge fails regardless of the calculations; also not considering the load path and building members that do not carry the load through the structure; at the highest level of competition, bridges often fail in torsion as the loading cannot be fully accounted for in the design \\
\addlinespace
\multicolumn{2}{@{}p{\linewidth}@{}}{\textbf{C4.} Does the club have models and recorded failure loads of past bridges that could be used to test the Toolbox?} \\*
S3-P1 & Some data available from the most recent couple of years; everything before about 3 years ago was never recorded by previous co-chairs/execs and is lost \\
S3-P2 & Some videos from past years and scoresheets from past competitions; most old bridges that have been kept have their ultimate load written on them \\
S3-P3 & Not really \\
\addlinespace
\multicolumn{2}{@{}p{\linewidth}@{}}{\textbf{CL1.} Is there anything about the club's information or tools that we did not ask about?} \\*
S3-P1 & N/A \\
S3-P2 & No \\
S3-P3 & \textit{No answer recorded} \\
\addlinespace
\multicolumn{2}{@{}p{\linewidth}@{}}{\textbf{CL2.} Are there any items that you feel are important to include in the Toolbox that we have not discussed?} \\*
S3-P1 & Long term: interactive structural calculation tools (input bridge configuration and design constraints to estimate member forces, deflections and potential failure points); team management features (milestone submissions, deadlines, coach feedback, progress tracking); searchable database of previous designs, material testing results and construction techniques; aware not all of this is feasible this year, hopes the software side of Troitsky develops over the next couple of years \\
S3-P2 & Checklists of analysis steps, understanding and rules may be useful for first-timers, though not strictly necessary \\
S3-P3 & \textit{No answer recorded} \\
\end{longtable}
